\mxnewpage\section*{Correction EV1-A}\label{EV1-A-CORR}
\textbf{QCM :} 1 B, 2 A, 3 D, 4 B, 5 A, 6 B, 7 B, 8 B, 9 D, 10 C, 11 A, 12 A. Chaque réponse correcte : \num{0,5} point ; total 6.
\paragraph{Vérifications des automatismes.} Priorités ; somme de fractions ; soustraction des exposants ; coefficient multiplicateur \num{1,2} ; division par \num{0,8} ; produit \num{1,25} fois \num{0,8} ; résolution affine ; accroissement vertical divisé par l’horizontal ; moyenne ; angle divisé par 360 ; complément à 1 ; substitution entre parenthèses.
\subsection*{Exercice 1 -- Réserve décroissante}
\begin{enumerate}
\item $u_0=192$, $u_{n+1}=\dfrac{u_n}{2}$ ; raison $\dfrac{1}{2}$. \textit{Barème : 1 point. Résultat et justification élémentaire attendus ; demi-point possible pour démarche correcte inachevée.}
\item $192,96,48,24,12$, somme 372. \textit{Barème : 2 points. Méthode ou justification : 1 ; calcul et conclusion : 1.}
\item $u_4=12$, $u_5=6$ : rang 5. Initialiser n=0,u=192 ; tant que u est au moins 10, diviser par 2 et augmenter n. \textit{Barème : 2 points. Méthode ou justification : 1 ; calcul et conclusion : 1.}
\end{enumerate}
\subsection*{Exercice 2 -- Paramètre et tangente}
\begin{enumerate}
\item $\Delta=4m^2-16=4(m^2-4)$. \textit{Barème : 1 point. Résultat et justification élémentaire attendus ; demi-point possible pour démarche correcte inachevée.}
\item Deux racines si $m<-2$ ou $m>2$ ; une si $m=\pm2$ ; aucune si $-2<m<2$. \textit{Barème : 2 points. Méthode ou justification : 1 ; calcul et conclusion : 1.}
\item Pour m=1, taux $2+h$, dérivée 2 et $f(2)=4$ ; tangente $y=2x$. \textit{Barème : 2 points. Méthode ou justification : 1 ; calcul et conclusion : 1.}
\end{enumerate}
\subsection*{Exercice 3 -- Locus par le milieu}
\begin{enumerate}
\item $MI^2-16$, par les identités scalaires. \textit{Barème : 1 point. Résultat et justification élémentaire attendus ; demi-point possible pour démarche correcte inachevée.}
\item $MI^2=25$ : cercle de centre I, rayon 5. \textit{Barème : 1 point. Résultat et justification élémentaire attendus ; demi-point possible pour démarche correcte inachevée.}
\item Pour -16, $MI^2=0$ : I seul. Pour -20, $MI^2=-4$ : vide. \textit{Barème : 2 points. Méthode ou justification : 1 ; calcul et conclusion : 1.}
\end{enumerate}
\paragraph{Étalonnage.} Durée cible : 20–25 min pour les automatismes, 80–85 min pour les exercices, 10–15 min de relecture. Ces durées sont des estimations éditoriales à mesurer en classe. Accepter toute méthode correcte ; ne pas sanctionner deux fois une erreur numérique propagée avec cohérence. Le barème proposé ne vaut pas décision d’harmonisation de l’équipe.
